\section{Application 2: A Smooth Divisor-Sum Function (\texorpdfstring{$\mathcal{P}_{\sigma}$}{P_sigma})}

\subsection{Motivation and Construction}
A smooth analogue of the sum-of-divisors function, $\sigma(n)$, can be constructed.
To achieve this, it is required that each divisor $d$ contributes its own value $d$ to the sum.
Since $F(n,d)=d^2$, the natural choice for the weighting function is $W(i)=1/i$.
\begin{definition}[Smooth Divisor-Sum Function]
For $x>0$ and a steepness parameter $\kappa>0$, define
\begin{equation}\label{eq:Psigma}
\mathcal{P}_{\sigma}(x; \kappa) = \left( \sum_{i=2}^{\infty} \phi_{\kappa}\!\left(\frac{i}{x+1}\right) \cdot \frac{F(x,i)}{i} \right) - x.
\end{equation}
\end{definition}
For integer arguments $n$, the summation part evaluates to the sum of all divisors except 1, which is $\sigma(n)-1$. Subtraction of $n$ yields an expression that vanishes precisely at primes, for which $\sigma(n)-1=n$.
\subsection{Properties of $\mathcal{P}_{\sigma}$}

\begin{proposition}[Properties of $\mathcal{P}_{\sigma}$]\label{prop:Psigma-properties}
The function $\mathcal{P}_{\sigma}(x;\kappa)$ is $C^\infty$ on $(0,\infty)$. For any integer $n\ge2$,
\[
\lim_{\kappa\to\infty} \mathcal{P}_{\sigma}(n; \kappa) = \sigma(n) - n - 1,
\]
and the right-hand side vanishes if and only if $n$ is prime.
\end{proposition}

\begin{proof}
The proof of $C^\infty$-regularity on $(0,\infty)$ follows the same argument as for $\mathcal{P}_{\tau}$. Let $h_i(x) = \phi_{\kappa}(i/(x+1)) F(x,i)/i$. For any compact $K\subset(0,\infty)$ and integer $j\ge0$, the Leibniz rule is applied. The bounds from Lemma~\ref{lem:cutoff-derivatives} and Proposition~\ref{prop:fejer-bounds} (where $|\partial_x^{j-q} (F(x,i)/i)| \le D_{j-q} \cdot i$) yield a uniform majorant on $K$:
\[
\sup_{x\in K} |\partial_x^j h_i(x)| \le M_{j,K} \, i^{j+1} \, e^{-c_K i}.
\]
This is summable for any fixed $j$. By the Weierstrass M-test, all derivative series converge uniformly on compact sets, establishing that $\mathcal{P}_{\sigma}$ is a $C^\infty$ function.

For an integer $n$, Proposition~\ref{prop:fejer-properties} implies $F(n,i)/i = i \cdot \mathbf{1}_{i\mid n}$. The series reduces to a finite sum over divisors:
\[
\mathcal{P}_{\sigma}(n;\kappa)=\sum_{\substack{d\mid n\\ d\ge2}} d\,\phi_{\kappa}\!\Bigl(\frac{d}{n+1}\Bigr)-n.
\]
In the limit $\kappa\to\infty$, $\phi_{\kappa}(d/(n+1))\to 1$ for all divisors $d$. The sum becomes $\sum_{d\mid n, d\ge2} d = \sigma(n)-1$. The entire expression thus converges to $\sigma(n)-n-1$.
\end{proof}

\begin{conjecture}[Asymmetric companion zeros for $\mathcal{P}_{\sigma}$]\label{conj:Psigma-companions}
For each odd prime $p\ge 3$ and sufficiently large $\kappa$, $\mathcal P_\sigma(\cdot;\kappa)$ displays an asymmetric pair of real zeros flanking $p$: an exponentially close left zero with
\[
p - x_p^-(\kappa)\ \asymp\ p\,\mathrm e^{-2\kappa/(p+1)}\qquad(\kappa\to\infty),
\]
whose existence and scale are proved in Lemma~\ref{lem:Psigma-left-zero}, and a right zero whose distance to $p$ appears to converge to a positive limit as $\kappa\to\infty$ (Figure~\ref{fig:PsigmaCompanions}). Partial results: for each fixed $p$ the right zero stays bounded away from $p$ (Proposition~\ref{prop:Psigma-right-fixed}) and is unique for $\kappa\ge3.25\,(p+1)\ln p$ (Proposition~\ref{prop:Psigma-unique}); the convergence of the distance and its uniformity in $p$ are open. The value at $p$ is explicit: $\mathcal P_\sigma(p;\kappa)=-p\bigl(1-\phi_\kappa(\tfrac{p}{p+1})\bigr)=-p/\bigl(1+\mathrm e^{2\kappa/(p+1)}\bigr)$.
\end{conjecture}

% remark: Unified terminology is provided once in Definition~\ref{def:zero-terminology}.

As a numerical illustration of Conjecture~\ref{conj:Psigma-companions}, 
Figure~\ref{fig:PsigmaCompanions} plots $\mathcal P_\sigma(x;\kappa)$ on $[2,8]$ for 
$\kappa\in\{2,10,20,100\}$. The predicted asymmetry is visible: an exponentially close left zero and a right zero at a distance bounded away from $0$ near each odd prime, while the value at $x=p$ remains negative for finite~$\kappa$. A complete scan of $[1.5,\,8.5]$ at $\kappa=40$ locates nine real zeros and shows that a zero lying to the left of a prime $p$ need not be unique: between $2$ and $3$ there are two, namely the right companion of $2$ at $x=2.70361851\ldots$ and the left approximant of $3$ at $x=3-6.18\cdot10^{-9}$. Accordingly, Lemma~\ref{lem:Psigma-left-zero} is an existence statement and $x_p^-(\kappa)$ denotes a zero produced by it, not a uniquely determined one.

\begin{figure}[!htbp]
\centering
\includegraphics[width=0.9\textwidth]{plot_companion_zeros_sigma.pdf}
\caption{Asymmetric companion zeros of $\mathcal{P}_{\sigma}(x;\kappa)$. Dashed lines mark primes $p=3,5,7$. Consistent with Conjecture~\ref{conj:Psigma-companions}, a left zero converges to $p$ exponentially fast with $\kappa$, while the right zero's distance to $p$ stabilizes at a positive limit. The value at $x=p$ is strictly negative for any finite $\kappa$.}
\label{fig:PsigmaCompanions}
\end{figure}

% remark: Quantitative statements are incorporated in Lemma~\ref{lem:uniform-local-sigma} and Lemma~\ref{lem:Psigma-left-zero}.

\begin{lemma}[Uniform local derivative bounds near an odd prime]\label{lem:uniform-local-sigma}
Fix a prime $p\ge 3$ and put $\varepsilon_0:=\tfrac12$. There exists $C_p>0$ such that, for all $\kappa\ge 1$ and all $x$ with $|x-p|\le \varepsilon_0$,
\[
\bigl|\partial_x^{2}\mathcal P_\sigma(x;\kappa)\bigr|\ \le\ C_p,\qquad
\bigl|\partial_x^{3}\mathcal P_\sigma(x;\kappa)\bigr|\ \le\ C_p\,\kappa .
\]
In particular, the quadratic Taylor coefficient $A_p(\kappa)=\tfrac12\,\mathcal P_\sigma''(p;\kappa)$ satisfies $|A_p(\kappa)|\le C_p/2$, and $|\mathcal P_\sigma(p+\varepsilon;\kappa)-\mathcal P_\sigma(p;\kappa)-\varepsilon\,\mathcal P_\sigma'(p;\kappa)|\le \tfrac12C_p\,\varepsilon^2$ for $|\varepsilon|\le\varepsilon_0$, uniformly in $\kappa\ge 1$. The third derivative is not bounded uniformly in $\kappa$: numerically, $\sup_{|x-5|\le0.3}|\partial_x^3\mathcal P_\sigma(x;\kappa)|$ grows from about $79$ at $\kappa=10$ to about $547$ at $\kappa=160$, while $\sup_{|x-5|\le0.3}|\partial_x^2\mathcal P_\sigma(x;\kappa)|$ stays between $22$ and $26$.
\end{lemma}

\begin{proof}
Put $\omega(y):=(1-\tanh y)/2$, so that $\phi_\kappa(u)=\omega(\kappa(u-1))$, write $u_i(x)=i/(x+1)$, $h_i(x):=\phi_\kappa(u_i(x))\,F(x,i)/i$, and let $K=[p-\tfrac12,p+\tfrac12]$ and $\kappa\ge1$. On $K$ one has $|u_i^{(k)}(x)|\le k!\,i$ for $k\ge1$, $|\omega^{(j)}(y)|\le c_j\,\mathrm e^{-2|y|}$ for $j\ge1$, and $|\partial_x^\ell(F(x,i)/i)|\le D_\ell\,i$ by Proposition~\ref{prop:fejer-bounds}. By Faà di Bruno's formula, $\partial_x^q\phi_\kappa(u_i(x))$ with $1\le q\le3$ is a sum of terms $\kappa^j\omega^{(j)}(\kappa(u_i-1))$, $1\le j\le q$, times products of derivatives of $u_i$. The series may be differentiated termwise (Proposition on $\mathcal P_\sigma$ above), and three groups of indices are distinguished.

\emph{(a) $2\le i\le p$.} Here $u_i(x)-1\le -\frac{1}{2p+1}$ on $K$, so $\kappa^j|\omega^{(j)}|\le c_j\kappa^j\mathrm e^{-2\kappa/(2p+1)}\le c_{j,p}$ uniformly in $\kappa$. Hence $|h_i''|+|h_i'''|\le C_p$ on $K$.

\emph{(b) $i\ge p+2$.} Here $v_i:=u_i(x)-1\ge\frac{i-p-3/2}{p+3/2}\ge\frac{1}{2p+3}$ on $K$, so $\phi_\kappa(u_i)\le\mathrm e^{-2\kappa v_i}$ and $\kappa^j|\omega^{(j)}|\le c_j\kappa^j\mathrm e^{-2\kappa v_i}\le c_j\,(j/v_i)^j\,\mathrm e^{-\kappa v_i}\le c_{j,p}\,\mathrm e^{-v_i}$ for $\kappa\ge1$. Hence $|h_i''|+|h_i'''|\le C_p\,i^{4}\,\mathrm e^{-(i-p-3/2)/(p+3/2)}$, which is summable in $i$ uniformly in $\kappa\ge1$.

\emph{(c) $i=p+1$.} With $t=x-p$ one has $s:=\kappa(u_{p+1}(x)-1)=-\kappa t/(x+1)$, hence $\kappa|t|\le(p+\tfrac32)|s|$. The function $G(x):=F(x,p+1)/(p+1)$ satisfies $G(p)=G'(p)=0$ (Proposition~\ref{prop:C1}), so $|G|\le C_pt^2$, $|G'|\le C_p|t|$ and $|G''|+|G'''|\le C_p$ on $K$. In the Leibniz expansion of $\partial_x^m(\omega(s)\,G)$ every term is bounded by $C_p\,\kappa^{j}|\omega^{(j)}(s)|\,|t|^{\max\{0,\,2-(m-q)\}}$ with $j\le q\le m$. For $m=2$ all these terms are $O_p(1)$, because $\sup_y|y|^k|\omega^{(j)}(y)|<\infty$; the largest one, $\kappa^2|\omega''(s)|\,t^2\le C_p\,s^2|\omega''(s)|$, is bounded. For $m=3$ the same estimate leaves one factor $\kappa$, e.g.\ $\kappa^3|\omega'''(s)|\,t^2\le C_p\,\kappa\,s^2|\omega'''(s)|$, so these terms are $O_p(\kappa)$.

Summing (a)--(c) and noting that the term $-x$ has vanishing second and third derivatives gives $|\partial_x^2\mathcal P_\sigma|\le C_p$ and $|\partial_x^3\mathcal P_\sigma|\le C_p\kappa$ on $K$.
\end{proof}


\begin{lemma}[Asymptotic derivative at primes]\label{lem:Psigma-derivative-at-p}
Fix a prime $p\ge 3$. The derivative of the smooth divisor-sum analogue satisfies
\[
\lim_{\kappa\to\infty} \mathcal{P}_\sigma'(p;\kappa) = -1.
\]
\end{lemma}
\begin{proof}
Let $S_\sigma(x;\kappa) = \sum_{i\ge 2}\phi_\kappa(i/(x+1))\,F(x,i)/i$. Then $\mathcal{P}_\sigma'(x;\kappa) = S_\sigma'(x;\kappa)-1$. The derivative of the sum is evaluated termwise at $x=p$. The derivative of the $i$-th term is
\[
\partial_x\Bigl(\phi_\kappa\!\Bigl(\frac{i}{x+1}\Bigr)\frac{F(x,i)}{i}\Bigr) =
\Bigl(\partial_x\phi_\kappa\!\Bigl(\frac{i}{x+1}\Bigr)\Bigr)\frac{F(x,i)}{i} + \phi_\kappa\!\Bigl(\frac{i}{x+1}\Bigr)\frac{F'(x,i)}{i}.
\]
At $x=p$, one has $F(p,i)=0$ for any $i\nmid p$ and $F'(p,i)=0$ for any $i$ (cf.\ Proposition~\ref{prop:C1}). Since $p$ is prime, the only divisor in the sum is $i=p$. Thus, only the term for $i=p$ could be non-zero. For $i=p$, $F(p,p)=p^2$ and $F'(p,p)=0$. The sum reduces to a single term:
\[
S_\sigma'(p;\kappa) = \left. \partial_x\phi_\kappa\!\left(\frac{p}{x+1}\right)\right|_{x=p} \cdot \frac{F(p,p)}{p}
= p \cdot \left(-\frac{p}{(p+1)^2}\right) \phi_\kappa'\!\left(\frac{p}{p+1}\right).
\]
Since $\phi_\kappa'(u) = -\kappa\,\mathrm{sech}^2(\kappa(u-1))/2$, one has $\phi_\kappa'(p/(p+1)) \asymp -\kappa e^{-2\kappa/(p+1)}$, which vanishes as $\kappa\to\infty$. Therefore, $\lim_{\kappa\to\infty} S_\sigma'(p;\kappa) = 0$, and the result follows.
\end{proof}


\begin{lemma}[Existence and scale of the left zero for $\mathcal P_\sigma$]\label{lem:Psigma-left-zero}
Let $p\ge3$ be prime. There exist $\kappa_0(p)\ge1$ such that for all $\kappa\ge\kappa_0(p)$ the function $\mathcal P_\sigma(\cdot;\kappa)$ has at least one real zero to the left of $p$; write $x_p^-(\kappa)$ for such a zero. Uniqueness among the zeros left of $p$ is not claimed. With
\[
\Delta_p^{(\sigma)}(\kappa):=p\Bigl(1-\phi_\kappa\!\Bigl(\tfrac{p}{p+1}\Bigr)\Bigr)\asymp p\,\mathrm e^{-2\kappa/(p+1)},
\]
the distance to $p$ is bounded by
\[
\frac{1}{2}\,\Delta_p^{(\sigma)}(\kappa)\ \le\ p-x_p^-(\kappa)\ \le\ 2\,\Delta_p^{(\sigma)}(\kappa).
\]
\end{lemma}

\begin{proof}
Set $\Delta:=\Delta_p^{(\sigma)}(\kappa)$. By Lemma~\ref{lem:uniform-local-sigma}, a Taylor expansion at $x=p$ is valid:
\[
\mathcal P_\sigma(p-\varepsilon;\kappa)=\mathcal P_\sigma(p;\kappa) - \varepsilon\,\mathcal P_\sigma'(p;\kappa) + O(\varepsilon^2),
\]
where the $O(\cdot)$ term is uniform in $\kappa\ge\kappa_0(p)$. One has $\mathcal P_\sigma(p;\kappa)=-\Delta$. By Lemma~\ref{lem:Psigma-derivative-at-p}, $\mathcal P_\sigma'(p;\kappa)=-1+o(1)$ as $\kappa\to\infty$. The expansion becomes
\[
\mathcal P_\sigma(p-\varepsilon;\kappa)=-\Delta + \varepsilon\,(1+o(1)) + O(\varepsilon^2).
\]
For $\varepsilon_1=\Delta/2$, the value is $\mathcal P_\sigma(p-\varepsilon_1;\kappa)=-\Delta/2+o(\Delta)<0$ for sufficiently large $\kappa$.
For $\varepsilon_2=2\Delta$, the value is $\mathcal P_\sigma(p-\varepsilon_2;\kappa)=\Delta+o(\Delta)>0$ for sufficiently large $\kappa$.
By the Intermediate Value Theorem, a zero exists in the interval $(p-2\Delta, p-\Delta/2)$.
\end{proof}

% remark: The status is already stated in Conjecture~\ref{conj:Psigma-companions}; no separate remark is needed.

\subsection{Critical discussion: non-integer zeros of $\mathcal{P}_{\sigma}$}
The indicator role of $\mathcal{P}_{\sigma}$ is confined to integer input. For non-integer $x$, the existence of zeros cannot be ruled out, because the function's sign is determined by the competition between the sum $S_{\sigma}(x;\kappa) = \sum_{i\ge2}\phi_{\kappa}(i/(x+1))\,F(x,i)/i$ and the linear term $x$. Near any integer, the factor $\sin^2(\pi x)$ in each $F(x,i)$ drives the sum to zero, ensuring $\mathcal{P}_{\sigma}(x;\kappa) \to -x$. Conversely, near rational points $x \approx im$ with small $i$, the term $F(x,i)/i$ can become large and lift $S_{\sigma}(x;\kappa)$ above the line $y=x$, creating intersections.

The following lemma provides a constructive lower bound for \emph{partial} sums of $S_{\sigma}$ away from such resonances. It does not yield a global sign for $\mathcal{P}_\sigma(x)=S_\sigma(x)-x$, as the two terms remain incomparable near integers and rational resonances. Consequently, no assertion about the sign of $\mathcal{P}_\sigma$ for non-integer arguments is claimed.

\begin{lemma}[Quantitative control away from resonances]\label{lem:constructive-lb}
Let $K=[a,b]\subset(0,\infty)$ be compact, $\eta\in(0,1)$, and $\kappa>0$. Choose $I_0\in\mathbb N$ and set $I:=\{2,\dots,I_0\}$. For $i\in I$ let
\[
E_i:=\{x\in K:\ \mathrm{dist}(x/i,\mathbb Z)<\delta_i\},\qquad 
\delta_i:=\frac{\eta}{8\,I_0\,\bigl((b-a)+I_0\bigr)}.
\]
Let $E_0:=\{x\in K:\ \mathrm{dist}(x,\mathbb Z)<\delta_0\}$ with $\delta_0:=\frac{\eta}{8(b-a+1)}$ and put $E:=E_0\cup\bigcup_{i\in I}E_i$. Then $\mathrm{meas}(E)\le \eta$ and, for all $x\in K\setminus E$,
\[
\sum_{i\in I}\phi_{\kappa}\!\Bigl(\frac{i}{x+1}\Bigr)\frac{F(x,i)}{i}
\ \ge\ 
\Bigl(\tfrac{2\delta_0}{\pi}\Bigr)^{\!2}\,
\sum_{i=2}^{I_0}\frac{\phi_{\kappa}\!\bigl(\tfrac{i}{a+1}\bigr)}{i}.
\]

\begin{proof}
For $x\in K\setminus E_0$, the distance $\mathrm{dist}(x,\mathbb Z)\ge \delta_0$ implies $\sin^2(\pi x)\ge (\tfrac{2\delta_0}{\pi})^2$. The map $x\mapsto \phi_\kappa(i/(x+1))$ is increasing on $[a,b]$ because $\phi_\kappa'(u)\le 0$ and $u'(x)<0$, hence
\[
\phi_\kappa\!\Bigl(\frac{i}{x+1}\Bigr)\ \ge\ \phi_\kappa\!\Bigl(\frac{i}{a+1}\Bigr)\qquad(x\in K).
\]
Using $F(x,i)=\sin^2(\pi x)/\sin^2(\pi x/i)$ and $\sin^2(\pi x/i)\le 1$ gives
\[
\phi_\kappa\!\Bigl(\frac{i}{x+1}\Bigr)\frac{F(x,i)}{i}
\ \ge\
\Bigl(\tfrac{2\delta_0}{\pi}\Bigr)^2\,\frac{\phi_\kappa\!\bigl(\tfrac{i}{a+1}\bigr)}{i}.
\]
Summing over $i\in I$ yields the lower bound. For the measure estimate, $\mathrm{meas}(E_0)\le 2\delta_0(b-a+1)$. For each $i\in I$, the condition $|x/i-m|<\delta_i$ is equivalent to $|x-im|<i\delta_i$, so $E_i$ is a union of at most $\lfloor(b-a)/i\rfloor+2$ intervals of length $2i\delta_i$. Therefore
\[
\mathrm{meas}\Bigl(\bigcup_{i=2}^{I_0}E_i\Bigr)\ \le\ \sum_{i=2}^{I_0}\Bigl(\Bigl\lfloor\frac{b-a}{i}\Bigr\rfloor+2\Bigr)\,2i\delta_i
\ \le\ 2\sum_{i=2}^{I_0}\bigl((b-a)+2i\bigr)\delta_i.
\]
With the chosen $\delta_i$, the sum is bounded by $\eta/2$, and together with $\mathrm{meas}(E_0)\le \eta/4$ this shows $\mathrm{meas}(E)\le \eta$.
\end{proof}
\end{lemma}

\begin{remark}[On the choice of bounds]
A lower bound for $\sin^2(\pi x)$ combined with the trivial upper bound $\sin^2(\pi x/i)\le 1$ yields a lower bound for $F(x,i)$. A lower bound on $\sin^2(\pi x/i)$ would only provide an upper bound for $F(x,i)$ and is therefore not useful here.
\end{remark}
\begin{remark}[On the choice of the thresholds $\delta_i$]
An $i$-dependent choice for $\delta_i$ could slightly improve the constants in the measure estimate. The uniform choice used above suffices for the proof and keeps the presentation simple. The sets $E_i$ are not needed for the lower bound itself, which only uses $x\notin E_0$; they are included so that the exceptional set also removes neighbourhoods of the resonances $x\approx im$.
\end{remark}

\begin{figure}[!htbp]
\centering
\includegraphics[width=0.9\textwidth]{plot_P_sigma_zeros.pdf}
\caption{Behavior of the smooth divisor-sum analogue $\mathcal{P}_{\sigma}(x; \kappa)$ for $\kappa=1000$ on $[0,8]$.
At any finite $\kappa$, $\mathcal P_\sigma(p;\kappa)<0$ at odd primes $p$; as $\kappa\to\infty$ the value at $x=p$ tends to $0$. Consistent with Conjecture~\ref{conj:Psigma-companions}, one observes a \emph{left prime-zero approximant} exponentially close to $p$ and a \emph{right companion zero} at a distance bounded away from zero. Away from these neighborhoods, the existence of additional non-integer zeros remains unresolved; accordingly, the prime-indicator interpretation is restricted to integer input.}
\label{fig:P_sigma}
\end{figure}

\FloatBarrier

% ---------------------------------------------------------------------------
